\subsection{Komponenten}
Komponenten erlauben die Speicherung und Modifikation domänenspezifischer Werte.\\

\noindent
Sei $K$ ein Komponentenschema.\\

\noindent
Stehe weiter mit $I$ eine Indexmenge von Bezeichnern (\textit{Labels}) zur Verfügung
Wir fassen damit $K_i$ als Familie auf, die über $i \in I$ indexiert ist.\\

\noindent
Bezeichne $k_i \in K_i$ einen beliebigen Term des Typs $K_i$.\footnote{``Each $K_i$ is a type in the sub-language, \ldots ``}
Redmond et al. definieren hierzu -- informell -- eine Abbildung:\footnote{
    ``The component types come to us in a mapping from component label $i \in I$ to component type $K_i$, \ldots ``
}

\[
    k: I \rightarrow K_i
\]

\noindent
wobei $K_i$ durch das Label $i \in I$ determiniert ist.
Wie in der zitierten Arbeit nutzen auch wir im Folgenden die Indexnotation $k_i$ anstelle $k(i)$:

\[
    k_i \in K_i
\]

\noindent
Hierdurch wird mit $k_i$ ein beliebiger, isolierter, zustandsloser Wert aus dem Wertebereich der durch den Komponententypen repräsentierten Domäne $K_i$ dargestellt.
Entsprechend kann ein beliebiger Term $k_i$ formal mit mehreren Entitäten assoziiert sein, ohne, dass es eines zusätzlichen Index bedarf.\footnote{
    In der Praxis wirkt sich dies als Kopie der entsprechenden Werte aus (Redmond et al.: ``Component values``).
}\\

\subsubsection*{Beispiel}
Redmond et al. notieren in ihrer Arbeit beispielhaft ein Komponentenschema $K$, das folgende Paare $i \mapsto K_i$ beinhaltet:
\[
    K \coloneqq \{ i_1 \mapsto K_{i_1}, i_2 \mapsto K_{i_2} \}
\]

\noindent
mit
\[
    \begin{align}
        i_1 \coloneqq \text{Pos}\\
        i_2 \coloneqq \text{Vel}\\
        K_{i_1}, K_{i_2} \coloneq \mathbb{Z}
    \end{align}
\]

\noindent
Für einen Komponententyp $K_\texttt{Position} \subseteq \mathbb{R}^2$ können wir also schreiben:

\[
    k_{\texttt{Position}} = (1.0, 1.0) \in K_\texttt{Position}
\]



\begin{tcolorbox}[title={Komponenten als Informationsträger für Referenzen}]
    Für Anwendungsfälle, in denen eine Komponente eine Referenz auf eine Objektinstanz $t$ vom Typ $T$ verwaltet, kann in diesem Kontext das Label

    \[
        \text{Ref(T)} \in I
    \]

    \noindent
    eingeführt werden.\\

    \noindent
    Hierzu sei

    \[
        \mathcal{A} \colon \text{Inst}(T) \rightarrow \mathbb{N}
    \]

    \noindent
    eine Abbildung, die eine Instanz eines Typs $T$ auf einen numerischen Wert abbildet -- den Speicherort der Instanz im \textit{Adressraum}.\\

    \noindent
    Ein Term $k_{\text{Ref(T)}}$ speichert dann diesen numerischen Wert für $t \in \text{Inst}(T)$

    \[
        k_{\text{Ref(T)}} = \mathcal{A}(t) \in K_{\text{Ref(T)}} = \mathbb{N}
    \]


    \noindent
    Dadurch ist die Frage, wie mit diesem Komponentenwert im ECS-Programm umgegangen wird, zunächst semantischer Natur, und unabhängig von der Tatsache, dass der Term tatsächlich einen Speicherort repräsentiert.
\end{tcolorbox}